\section{Optimasi dan VQE}

% Slide 9.1: Bab 2.6.1 - Gradien Energi
\begin{frame}{Bab 2.6.1 - Gradien Energi}
    \begin{itemize}
        \item Pembaruan parameter menggunakan nilai dari gradien energi
        \begin{equation}
            \frac{\partial E}{\partial \theta} = \frac{\langle\psi(\theta)|\hat{H}|\psi(\theta)\rangle}{\partial \theta}
        \end{equation}
        \item Misal sistem memliki 2 qubit, maka 
        \begin{equation}
            \begin{split}
            \frac{\partial E}{\partial \theta} &= \frac{\bra{00} \hat{U}_{A}^\dagger \hat{U}(\theta)^\dagger \hat{U}_{B}^\dagger \hat{\mathcal{H}} \hat{U}_{B} \hat{U}(\theta) \hat{U}_{A} \ket{00}}{\partial \theta} \\
            &= \frac{\bra{\phi} \hat{U}(\theta)^\dagger \hat{U}_{B}^\dagger \hat{\mathcal{H}} \hat{U}_{B} \hat{U}(\theta) \ket{\phi}}{\partial \theta} \\
            &= \frac{\bra{\phi} \hat{U}(\theta)^\dagger \hat{M} \hat{U}(\theta) \ket{\phi}}{\partial \theta}
            \end{split}
        \end{equation}
    \end{itemize}
\end{frame}

\begin{frame}{Bab 2.6.1 - Gradien Energi Lanjutan}
    \begin{itemize}
        \item dengan turunan uniter
        \begin{equation*}
            \frac{\partial U(\theta)}{\partial\theta} = -\frac{i}{2}\sigma(\theta) U(\theta) \quad \text{dan} \quad \frac{\partial U^\dagger(\theta)}{\partial\theta} = \frac{i}{2}\sigma(\theta) U^\dagger(\theta)
        \end{equation*}
    \item maka turunan menjadi 
    {\tiny
    \begin{equation*}
        \begin{split}
        \frac{\partial \langle E(\theta)\rangle}{\partial \theta} &= \bra{\phi} \frac{\partial U^{\dagger}(\theta)}{\partial \theta} M U(\theta) \ket{\phi} + \bra{\phi} U^{\dagger}(\theta) M \frac{\partial U(\theta)}{\partial \theta} \ket{\phi} \\
        &= \bra{\phi} \left(\frac{i}{2}\sigma U^{\dagger}(\theta)\right) M U(\theta) \ket{\phi} + \bra{\phi} U^{\dagger}(\theta) M \left(-\frac{i}{2}\sigma U(\theta)\right) \ket{\phi} \\
        &= \frac{i}{2} \bra{\psi(\theta)} \sigma M \ket{\psi(\theta)} - \frac{i}{2} \bra{\psi(\theta)} M \sigma \ket{\psi(\theta)} \\
            &= \frac{i}{2} \bra{\psi(\theta)} \sigma M - M \sigma \ket{\psi(\theta)}
        \end{split}
    \end{equation*}
    }
    \item sehingga didapatkan
        \begin{equation}
        \frac{\partial \langle E(\theta)\rangle}{\partial \theta} = \frac{i}{2} \bra{\psi(\theta)} [\sigma, M] \ket{\psi(\theta)}
        \end{equation}
\end{itemize}
\end{frame}

% Slide 9.2: Bab 2.6.2 - Parameter Shift Rule (PSR)
\begin{frame}{Bab 2.6.2 - Parameter Shift Rule (PSR)}
    \begin{itemize}
        \item PSR menggunakan konsep pergeseran parameter $\theta$ sebesar $s = \frac{\pi}{2}$ ke arah positif dan negatif
        \begin{equation}
            \begin{split}
            \hat{U}(\theta \pm s) &= \exp\left(-i\frac{\theta \pm s}{2}\sigma\right) \\
            &= \exp\left(-i\frac{\theta}{2}\sigma\right) \exp\left(\mp i\frac{s}{2}\sigma\right) \\
            &= \hat{U}(\theta) \hat{U}(\pm s).
            \end{split}
        \end{equation}
        \item sehingga ekspektasi energi menjadi
        \begin{equation*}
            \begin{split}
            \langle E(\theta\pm s)\rangle &= \bra{\phi}\hat{U}(\theta\pm s)^{\dagger}\hat{M}\hat{U}(\theta\pm s)|\phi\rangle \\
            &= \bra{\phi}\hat{U}(\theta)^{\dagger}\hat{U}(\pm s)^{\dagger}\hat{M}\hat{U}(\theta)\hat{U}(\pm s)|\phi\rangle \\
            &= \bra{\psi} \hat{U}(\pm s)^{\dagger}\hat{M}\hat{U}(\pm s)|\psi\rangle.
            \end{split}
        \end{equation*}
    \end{itemize}
\end{frame}

\begin{frame}{Bab 2.6.2 - Parameter Shift Rule (PSR)}
    \begin{itemize}
        \item sehingga untuk pergeseran positif
        {\tiny
        \begin{equation}
            \begin{split}
                \langle E(\theta+s)\rangle = \bra{\psi(\theta)} & \left[ \cos\left(\frac{s}{2}\right) \hat{I} + i \sin\left(\frac{s}{2}\right) \sigma \right] \hat{M} \left[ \cos\left(\frac{s}{2}\right) \hat{I} - i \sin\left(\frac{s}{2}\right) \sigma \right] \ket{\psi(\theta)} \\
                = \bra{\psi(\theta)} & \left[ \cos^2\left(\frac{s}{2}\right) \hat{M} - i \cos\left(\frac{s}{2}\right) \sin\left(\frac{s}{2}\right) \hat{M} \sigma + \right. \\
                & \left. i \sin\left(\frac{s}{2}\right) \cos\left(\frac{s}{2}\right) \sigma \hat{M} + \sin^2\left(\frac{s}{2}\right) \sigma \hat{M} \sigma \right] \ket{\psi(\theta)}.
            \end{split}
        \end{equation}
        \begin{equation}
            \begin{split}
            \langle E(\theta+s)\rangle = &\cos^2\left(\frac{s}{2}\right) \langle \hat{M} \rangle + \sin^2\left(\frac{s}{2}\right) \langle\sigma \hat{M} \sigma\rangle + i \sin\left(\frac{s}{2}\right) \cos\left(\frac{s}{2}\right) \langle [\sigma, \hat{M}] \rangle.
            \end{split}
        \end{equation}
        \begin{equation}
            \begin{split}
            \langle E(\theta-s)\rangle = &\cos^2\left(\frac{s}{2}\right) \langle \hat{M} \rangle + \sin^2\left(\frac{s}{2}\right) \langle\sigma \hat{M} \sigma\rangle - i \sin\left(\frac{s}{2}\right) \cos\left(\frac{s}{2}\right) \langle [\sigma, \hat{M}] \rangle.
            \end{split}
        \end{equation}
        }
        \item dengan menghitung selisih kedua pergeserang tersebut, maka didapatkan;
        \begin{equation}
            \begin{split}
                \langle E(\theta+s)\rangle - \langle E(\theta-s)\rangle &= 2 i \sin\left(\frac{s}{2}\right) \cos\left(\frac{s}{2}\right) \langle [\sigma, \hat{M}] \rangle \\
                &= i \sin(s) \langle [\sigma, \hat{M}] \rangle.
            \end{split}
        \end{equation}
    \end{itemize}
\end{frame}

\begin{frame}{Bab 2.6.2 - Parameter Shift Rule (PSR)}
    \begin{itemize}
        \item Jika $s = \frac{\pi}{2}$, maka $\sin(s) = 1$ dan $\cos(s) = 0$, sehingga didapatkan persamaan umum parameter shift rule
        \begin{equation*}
            \begin{split}
                2 \frac{\partial \langle E(\theta)\rangle}{\partial \theta} &= i \langle [\sigma, M]\rangle \\
                2 \frac{\partial \langle E(\theta)\rangle}{\partial \theta} &= \langle E(\theta+s)\rangle - \langle E(\theta-s)\rangle \\
            \end{split}
        \end{equation*}
        \begin{equation}
            \frac{\partial \langle E(\theta)\rangle}{\partial \theta} = \frac{1}{2} \left[ \langle E(\theta+s)\rangle - \langle E(\theta-s)\rangle \right]
        \end{equation}
    \end{itemize}
\end{frame}

% Slide 9.3: Bab 2.6.3 - Gradient Descent (GD)
\begin{frame}{Bab 2.6.3 - Gradient Descent (GD)}
    \begin{itemize}
        \item \textbf{Mekanisme:} Parameter $\theta$ diperbarui searah negatif gradien untuk mencapai minimum energi.
        \begin{equation*}
            \theta_{t+1} = \theta_t - \eta \nabla E(\theta_t)
        \end{equation*}
        \item dimana 
        \begin{itemize}
            \item $\theta_{t+1}$ adalah parameter pada iterasi ke $t+1$
            \item $\eta$ adalah laju pembelajaran 
            \item $\nabla E(\theta_t)$ adalah gradien energi pada parameter $\theta_t$.
        \end{itemize}
        \item \textbf{Kekurangan:} Harus mengevaluasi gradien dua kali untuk setiap parameter, sehingga kompleksitas evaluasi gradien PSR tumbuh linear terhadap jumlah parameter.
    \end{itemize}
\end{frame}

% Slide 9.4: Bab 2.6.4 - Optimizer SPSA
\begin{frame}{Bab 2.6.4 - Optimizer SPSA}
    \begin{itemize}
        \item \textbf{Inovasi:} Estimasi gradien seluruh parameter secara simultan menggunakan vektor perturbasi acak ($\pm 1$).
        {\tiny
        \begin{equation}
            \begin{split}
                E_{\pm} &= E(\theta_k \pm c_k \Delta_k) \\ 
                &\approx E(\theta_k) \pm c_k \sum_{j=1}^p \Delta_j \frac{\partial E}{\partial \theta_j} + \frac{1}{2} c_k^2 \sum_{j=1}^p \Delta_j^2 \frac{\partial^2 E}{\partial \theta_j^2} + \dots
            \end{split}
        \end{equation}
        }
        \item Sehingga jika diambil selisih kedua evaluasi tersebut,
        {\tiny
        \begin{equation}
            \begin{split}
                E_+ - E_- &= 2 c_k \sum_{j=1}^p \Delta_j \frac{\partial E}{\partial \theta_j} + O(c_k^3)
            \end{split}
        \end{equation}
        }
        \item maka didapatkan bentuk estimator gradien
        {\tiny
        \begin{equation}
            \begin{split}
                \hat{g}_i &= \frac{E_+ - E_-}{2 c_k \Delta_i} \\
                &= \frac{1}{\Delta_i} \sum_{j=1}^p \Delta_j \frac{\partial E}{\partial \theta_j} \\
                &= \frac{\Delta_i}{\Delta_i} \frac{\partial E}{\partial \theta_i} + \sum_{j \neq i} \frac{\Delta_j}{\Delta_i} \frac{\partial E}{\partial \theta_j} \\
                \mathbb{E} \left[ \hat{g}_i \right] &= \frac{\partial E}{\partial \theta_i} + \sum_{j \neq i} \mathbb{E} \left[ \frac{\Delta_j}{\Delta_i} \right] \frac{\partial E}{\partial \theta_j}
            \end{split}
        \end{equation}
        }
    \end{itemize}
\end{frame}

\begin{frame}{Bab 2.6.4 - Optimizer SPSA}
    \begin{itemize}
        \item Sehingga dekspektasi estimasi gradien
        \begin{equation}
            \mathbb{E} \left[ \hat{g}_i \right] = \frac{\partial E}{\partial \theta_i}
        \end{equation} 
        \item  dengan proses optimasi 
        \begin{equation}
            \theta_{k+1} = \theta_k - \eta \hat{g}_k
        \end{equation}
        \item dengan 
        \begin{itemize}
            \item $c_k$ merupakan koefisien perturbasi yang menurun seiring dengan iterasi 
            \item $\eta$ merupakan laju pembelajaran 
            \item $\hat{g}_k$ merupakan estimasi gradien pada iterasi ke $k$
        \end{itemize}
        \item SPSA lebih efisien dibandingkan GD karena hanya membutuhkan 50\% evaluasi fungsi biaya pada setiap iterasi dari GD
    \end{itemize}
\end{frame}

% Slide 10.2: Bab 2.7 - Arsitektur EfficientSU2 (HEA)
\begin{frame}{Bab 2.7 - VQE dengan Arsitektur EfficientSU2 (HEA)}
    \begin{itemize}
        \item Menggunakan prinsip variasi Reyleigh-Ritz yang menyatakan bahwa energi dari status percobaan selalu lebih besar atau sama dengan energi \textit{ground state} sejati ($E_0$).
        \begin{equation*}
            E(\theta) = \frac{\langle\psi(\theta)|\hat{H}|\psi(\theta)\rangle}{\langle\psi(\theta)|\psi(\theta)\rangle} \ge E_0
        \end{equation*}
        \item Menggunakan struktur lapisan rotasi ($R_y, R_z$) dan lapisan keterbelitan (\textit{entanglement}) menggunakan gerbang CNOT.
    \end{itemize}
    \begin{center}
        \includegraphics[width=0.6\textwidth]{../GTQuantumInvest/Hasil_N2_GT/Analisis_Window_N2/2021-05-05_circuit.png}
    \end{center}
\end{frame}

% Slide 11.1: Bab 2.8.1 - Fenomena Barren Plateau
\begin{frame}{Bab 2.8.1 - Fenomena Barren Plateau dan Keacakan}
    \begin{itemize}
        \item Semakin tinggi \textbf{kedalaman sirkuit} = semakin besar \textbf{ekspresibilitas} = semakin besar risiko terjebak pada \textbf{Barren Plateau}
        \item \textbf{Barren Plateau} ditandai oleh lenyapnya varians gradien karena  
        \begin{equation}
            \text{Var}[\partial_{\theta} E] \propto \frac{1}{2^{NL}}
        \end{equation}
        dimana N adalah jumlah qubit dan L adalah kedalaman sirkuit
        \item Untuk mengkuran keacakan keadaaan kuantum selama optimasi, digunakan entropi von Neumann
        \begin{equation}
            S(\rho) = -\text{Tr}(\rho \log \rho)
        \end{equation}
        dimana $\rho$ merupakan matriks densitas dari keadaan kuantum
    \end{itemize}
\end{frame}

% Slide 11.2: Bab 2.8.2 - Diagnostik: Entropi Von Neumann
\begin{frame}{Bab 2.8.2 - Entropi Von Neumann dan Entropi Shannon}
    \begin{itemize}
        \item Mengukur derajat kemurnian state kuantum selama optimasi
        \item Berawal dari entopi Shannon, 
    \end{itemize}
    {\scriptsize
    \begin{equation*}
        \begin{split}
            H &= - \sum_i p_i \log_2 (p_i) \\
            \rho &= \sum_i p_i |\psi_i\rangle \langle \psi_i| \\
            \rho &= \sum_j \lambda_j |e_j\rangle \langle e_j| \\
            \log_2(\rho) &= \sum_j \log_2 (\lambda_j) |e_j\rangle \langle e_j| \\
            \rho \log_2(\rho) &= \sum_j \lambda_j \log_2(\lambda_j) |e_j\rangle \langle e_j| \\
            \text{Tr}(\rho \log_2(\rho)) &= \sum_j \lambda_j \log_2(\lambda_j) \\
        \end{split}
    \end{equation*}
    \begin{equation}
        \therefore S(\rho) = -\text{Tr}(\rho \log_2(\rho)) = -\sum_j \lambda_j \log_2 (\lambda_j)
    \end{equation}
    }
\end{frame}
